% HISTORICAL author-side fragment; superseded by the integrated and checked
% parts/supports/supports_transfer.tex, which the master actually includes.
% AI-produced by GPT-6 Astra (OpenAI), ChatGPT Pro mode, 19 November 2026.
% Additional Part A proof. No priority claim. Uses the master's theorem counter.

\subsection{The split-poset comparison and the finite classification}
\label{sec:split-poset-comparison}

For a poset $P$, write
\[
 Q^+(P)=\{[u<v]:u,v\in P,\ u<v\},
 \qquad [u<v]\perp[r<s]\iff v\le r\ \text{or}\ s\le u.
\]
This is the orthogonality space of proper quotients in
\citet[\S4, p.~6]{Jenca2022}. Its order by interval inclusion is a
different relation. We write $\mathcal C(O)$ for the complete
ortholattice of double-polar closed subsets of an orthogonality space
$O$; complements and polars below are always taken in the indicated
ambient space.

\begin{theorem}[Proved: the comparison for every ordered set]
\label{SA21-MT01}
For an arbitrary poset $L$, let $S(L)=L\times\{-,+\}$ with
$x^-<x^+$ and $x^\epsilon<y^\delta$ whenever $x<y$ in $L$.
The maps
\[
 P_x\longmapsto[x^-<x^+],\qquad
 I_{a,b}\longmapsto[a^+<b^-]\quad(a<b)
\]
embed the disjoncteur of $M(L)$ as an induced subspace $U$ of
$Q^+(S(L))$. They induce an isomorphism of complete ortholattices
\[
 \Sigma_{M(L)}\cong\mathcal C\bigl(Q^+(S(L))\bigr).
\]
The isomorphism and its inverse are given explicitly by trace and by
the finite-piece extension described in the proof. No finiteness,
boundedness, connectedness or existence of meets and joins is assumed.
\end{theorem}
\begin{proof}
The three types of orthogonality in $M(L)$ are
\begin{align*}
 P_x\perp P_y&\iff x<y\ \text{or}\ y<x,\\
 P_x\perp I_{a,b}&\iff x\le a\ \text{or}\ b\le x,\\
 I_{a,b}\perp I_{c,d}&\iff b\le c\ \text{or}\ d\le a.
\end{align*}
Substitution in the displayed quotient relation gives exactly these
conditions, including the endpoint equalities. Thus the embedding is
an induced orthogonality embedding.

Every $q=[a^\epsilon<b^\delta]$ with $a<b$ is a concatenation of
at most three members of $U$: include $[a^-<a^+]$ if $\epsilon=-$,
always include $[a^+<b^-]$, and include $[b^-<b^+]$ if $\delta=+$.
Call this finite set $D(q)$. If both endpoints have the same base
element, then $q=[a^-<a^+]\in U$ and $D(q)=\{q\}$.
In every case the pieces have successive common endpoints and together
run from the first endpoint of $q$ to its last endpoint. Consequently,
for every $r\in Q^+(S(L))$,
\begin{equation}
 \label{eq:MT-pieces}
 q\perp r\quad\Longleftrightarrow\quad
 d\perp r\text{ for every }d\in D(q).
\end{equation}
The forward implication follows from interval containment. For the
reverse implication, write the pieces as
$[c_0<c_1],\ldots,[c_{k-1}<c_k]$ and $r=[v<w]$.
For each piece either $c_i\le v$ or $w\le c_{i-1}$.
If both alternatives occur at consecutive positions, they imply
$w\le c_i\le v$, contradicting $v<w$. Thus all pieces lie on the
same side of $r$, giving $c_k\le v$ or $w\le c_0$, as required.

For $B\subseteq U$, set
\[
 E(B)=\{q\in Q^+(S(L)):D(q)\subseteq B\}.
\]
We identify $U$ with $M(L)$ via the embedding. Equation
\eqref{eq:MT-pieces} gives
$B^{\perp_{Q}}=E(B^{\perp_U})$.
More generally,
\[
 E(B)^{\perp_Q}=E(B^{\perp_U}).
\]
Indeed testing orthogonality to $E(B)$ includes testing every $b\in B$
because $D(b)=\{b\}$; conversely the pieces of every element of
$E(B)$ belong to $B$, so the same test suffices. Hence
$E(B)^{\perp_Q\perp_Q}=E(B^{\perp_U\perp_U})$.
In particular $E$ takes closed sets to closed sets, commutes with
complement, and its trace on $U$ is $B$.

Conversely let $A\subseteq Q^+(S(L))$ be closed. If $q\in A$, then
every piece of $q$ is orthogonal to $A^{\perp_Q}$ by
\eqref{eq:MT-pieces}, and hence belongs to $A$. Conversely if every
piece lies in $A$, that same equation gives $q\in A$. Therefore
$A=E(A\cap U)$. The preceding polar formula then shows that $A\cap U$
is closed in $U$. Thus trace and $E$ are mutually inverse order
isomorphisms preserving orthocomplement; consequently they preserve
all meets and joins. This proves the complete-ortholattice comparison,
including the empty case.
\end{proof}

We call $P$ a \emph{conditional lattice} when every pair having a common
upper bound has a least upper bound and every pair having a common
lower bound has a greatest lower bound. We do not require a global
least or greatest element.

\begin{lemma}[Proved: finite conditional lattices suffice for quotient spaces]
\label{SA21-MT02}
If $P$ is a finite conditional lattice, then $\mathcal C(Q^+(P))$
is orthomodular.
\end{lemma}
\begin{proof}
Adjoin two new strict bounds $\widehat0,\widehat1$ to form
$\widehat P$. This is a finite lattice: joins or meets already present
in $P$ are retained, and a pair with no upper, respectively lower,
bound in $P$ has join $\widehat1$, respectively meet $\widehat0$.
Let $X$ be closed in $Q^+(P)$ and put
$\widehat X=X^{\perp\perp}$ in $Q^+(\widehat P)$.
Then
\begin{equation}
 \label{eq:MT-trace}
 \widehat X\cap Q^+(P)=X.
\end{equation}
To verify the possibly delicate inclusion, a point of the left side
is orthogonal to every element of $X^{\perp}$ in the larger space,
hence to every element of its trace, which is exactly the polar in
$Q^+(P)$. Since $X$ is closed in the smaller space, the point belongs
to $X$. The reverse inclusion is extensivity of closure.

By \citet[Theorem~4.10(a)$\Rightarrow$(b), p.~10]{Jenca2022}, the
closed set $\widehat X$ is of chain type: it is a finite union of
full quotient intervals
\[
 [a_1<b_1]^{\downarrow}\ \cup\cdots\cup\
 [a_k<b_k]^{\downarrow},
 \qquad a_1<b_1<a_2<b_2<\cdots<a_k<b_k
\]
in $\widehat P$. We have retained both finiteness and boundedness
when invoking this theorem. Taking traces gives an orthogonal union
$X=X_1\cup\cdots\cup X_k$. A nonempty trace block is one of the
following: all quotients in $[a,b]\subseteq P$; all quotients whose
upper endpoint is at most $b$; all quotients whose lower endpoint is
at least $a$; or all of $Q^+(P)$. The last three cases arise only from
a new endpoint $\widehat0$ or $\widehat1$.

Let $B$ be any maximal pairwise orthogonal subset of $X$. Since the
blocks are mutually orthogonal, $B_i=B\cap X_i$ is maximal in $X_i$.
We show directly that $B_i^{\perp}=X_i^{\perp}$, with both polars
taken in $Q^+(P)$; we do not assume that an individual trace block
is itself closed. The case $X_i=\varnothing$ is immediate. In a
nonempty block, order the finite quotient family $B_i$ from left to
right. Its intervals have successive endpoints
\[
 [c_0<c_1],\ [c_1<c_2],\ \ldots,\ [c_{r-1}<c_r].
\]
There can be no gap between consecutive intervals, since the quotient
across that gap would lie in $X_i$ and enlarge $B_i$.
For a bounded block $[a,b]$, maximality also forces $c_0=a,c_r=b$.
By the concatenation argument in \eqref{eq:MT-pieces}, a quotient is
orthogonal to $B_i$ precisely when it lies below $a$ or above $b$;
that is exactly orthogonality to the entire block.

For a lower-ray block, maximality forces $c_r=b$ and makes $c_0$
minimal in all of $P$: any $p<c_0$ would still satisfy $p<b$ and
would supply an additional quotient in this block. Thus no proper
quotient of $P$ can lie below $c_0$. Its polar consists exactly of
the quotients above $b$, the polar of the whole block. The upper-ray
case is dual. For the whole space, a maximal orthogonal family has
no common orthogonal element, so both polars are empty.

We now intersect these equalities:
\[
 B^{\perp}=\bigcap_i B_i^{\perp}
             =\bigcap_i X_i^{\perp}=X^{\perp}.
\]
As $X$ is closed, $B^{\perp\perp}=X$. The Dacey criterion
\citep[Theorem~3.3, p.~5]{Jenca2022} therefore gives orthomodularity.
This argument also treats disconnected $P$, empty blocks and the
empty poset. It does not infer heredity of orthomodularity for an
arbitrary induced subspace.
\end{proof}

\begin{theorem}[Proved: the exact finite ordered-set class]
\label{SA21-MT03}
For every finite poset $L$,
\[
 \Sigma_{M(L)}\text{ is orthomodular}
 \quad\Longleftrightarrow\quad
 L\text{ is a conditional lattice}.
\]
The conclusion does not require $L$ to be bounded or connected.
\end{theorem}
\begin{proof}
First, $S(L)$ is a conditional lattice exactly when $L$ is one.
For copies of comparable base elements, their meet and join are their
minimum and maximum. For incomparable $x,y$ with a common lower
bound, the meet of any two copies is $(x\wedge y)^+$; dually, with a
common upper bound the join is $(x\vee y)^-$. These formulas prove
sufficiency of the condition on $L$. Conversely, projecting a meet or
join of two copies to its base element gives the corresponding
greatest lower or least upper bound in $L$. Consequently the forward
construction in Lemma~\ref{SA21-MT02}, followed by
Theorem~\ref{SA21-MT01}, proves the required sufficient direction.

For necessity, suppose first that a pair $x,y$ has common upper
bounds but no least upper bound. Write
\[
 U=\{x,y\}^{\uparrow},\qquad I=U^{\downarrow}.
\]
Then $I^{\uparrow}=U$, because $x,y\in I$, and $U^{\downarrow}=I$.
The lower ideal $I$ has no greatest element; otherwise that element
would be the missing join of $x,y$. Similarly $U$ has no least
element, and $I\cap U=\varnothing$. Since $I$ is finite and nonempty,
choose two distinct maximal elements $a,b\in I$. They are incomparable.
Choose a minimal element $m$ of $L$ below $a$; it belongs to $I$.

Inside $Q^+(S(L))$, let
\[
 X=Q^+(S(I)),\qquad Y=Q^+(S(U)).
\]
We claim that $X^{\perp}=Y$ and $Y^{\perp}=X$. Every base element
of $I$ is strictly below every base element of $U$, so $X\perp Y$.
Conversely suppose $q$ is orthogonal to $X$. For every $i\in I$
it is orthogonal to the vertical quotient $[i^-<i^+]$.
It cannot lie below any such quotient: then both of its endpoint
bases would lie in the lower ideal $I$, making $q\in X$, contrary
to $q\perp X$. Therefore $q$ lies above every vertical quotient of
$I$, and its lower endpoint base belongs to $I^{\uparrow}=U$.
As $U$ is an upper ideal, both endpoint bases belong to $U$, so
$q\in Y$. The dual reasoning gives $Y^{\perp}=X$. Hence $X$ is
orthoclosed.

The singleton
\[
 B=\{[m^-<a^+]\}
\]
is a maximal pairwise orthogonal subset of $X$. A further quotient
orthogonal to it would have to lie below $m^-$ or above $a^+$.
The first is impossible because $m$ is minimal in $L$, and the
second is impossible within $S(I)$ because $a$ is maximal in $I$.
Choose any $u\in U$. The quotient $[a^+<u^-]$ is orthogonal to
$B$, but it is not orthogonal to $[b^-<b^+]$: the alternatives would
require $u\le b$ or $b\le a$, both false. Thus
\[
 [b^-<b^+]\in X\setminus B^{\perp\perp}.
\]
The Dacey criterion fails, so the quotient logic, and therefore
$\Sigma_{M(L)}$, is not orthomodular. If instead a lower-bounded
pair lacks a meet, apply the same argument to the opposite order;
reversing quotient endpoints gives an isomorphism of the two
orthogonality spaces. This proves necessity and the equivalence.
\end{proof}

The finite certificate program additionally constructs this cut witness
whenever it detects a missing bounded join (or dually meet), and verifies
the polar, maximality and non-spanning assertions. The result above is
a theorem for every finite cardinality, not an extrapolation from the
earlier enumeration.

\begin{remark}[Source hypothesis audit]
\label{SA21-MT04}
We use the bounded theorem of Jen\v ca, not the unrestricted wording
of his intermediate Corollary~4.5. As printed on p.~8, that corollary
asserts $q^{\perp\perp}=q^{\downarrow}$ for every poset. Two disjoint
two-element chains provide a finite counterexample: their quotient
space has two elements and no orthogonality, so the double polar of
either singleton is both elements, whereas its interval ideal is a
singleton. The preceding Lemma~4.3 and Corollary~4.4 explicitly impose
the relevant boundedness. Our transfer and trace-block arguments do
not use the invalid unbounded extension. No assertion here extends
the finite conditional-lattice characterization to arbitrary infinite
posets; Theorem~\ref{SA21-MT01} itself is unrestricted.
\end{remark}
